% B03-S001 | APPM5720Notes.tex baris 2971-2972
% segment-id: d90.becker.98ed693.b03.seg0001
\chapter{Reduksi Varians untuk SAA}
\label{becker03:chapter:variance-reduction}

Pertimbangkan masalah jumlah hingga dari aproksimasi rata-rata sampel
(\emph{sample average approximation}, SAA)
\begin{equation}
  \label{becker03:eq:finite-sum}
  \min_{x\in\R^n} f(x),
  \qquad
  f(x)=\frac{1}{N}\sum_{i=1}^{N}f_i(x),
\end{equation}
dengan $f$ bernilai hingga pada domain model. Sebagai contoh, untuk model
dengan prediktor linear dan fungsi kerugian dapat dipakai
\begin{equation}
  \label{becker03:eq:linearized-loss}
  f_i(x)=\ell(a_i^\top x-b_i),
\end{equation}
dengan $\ell$ fungsi kerugian. Huruf $\ell$ menggantikan huruf $L$ pada donor
agar tidak bertabrakan dengan konstanta kemulusan Lipschitz yang dipakai di
bawah.

% B03-S002 | APPM5720Notes.tex baris 2974-2981
% segment-id: d90.becker.98ed693.b03.seg0002
\section{Tabel gradien dan penaksir SAGA}
\label{becker03:section:saga-estimator}

Pilih iterat awal $x_0$ dan inisialisasikan satu titik tabel
$\phi_i^0=x_0$ untuk setiap $i=1,\dots,N$. Simpan, pada iterasi $k\geq0$,
\begin{equation}
  g_i^k=\nabla f_i(\phi_i^k),
  \qquad
  \bar g_k=\frac{1}{N}\sum_{i=1}^{N}g_i^k.
\end{equation}
Dengan $x\in\R^n$, tabel lengkap mempunyai $N$ kolom gradien di $\R^n$.
Ambil indeks acak
$J_k\sim\operatorname{Unif}\{1,\dots,N\}$ dan bentuk
\begin{equation}
  \label{becker03:eq:saga-estimator}
  v_k
  =\nabla f_{J_k}(x_k)-g_{J_k}^k+\bar g_k.
\end{equation}
Pembaruan SAGA tanpa suku proksimal adalah
\begin{equation}
  \label{becker03:eq:saga-update}
  x_{k+1}=x_k-t_kv_k.
\end{equation}
Sesudah $v_k$ dihitung dari tabel lama, hanya kolom terpilih yang disegarkan:
\begin{equation}
  \label{becker03:eq:table-update}
  g_i^{k+1}=
  \begin{cases}
    \nabla f_i(x_k),&i=J_k,\\
    g_i^k,&i\ne J_k,
  \end{cases}
  \qquad
  \phi_i^{k+1}=
  \begin{cases}
    x_k,&i=J_k,\\
    \phi_i^k,&i\ne J_k.
  \end{cases}
\end{equation}
Ini menuliskan secara eksplisit inisialisasi dan urutan pembaruan yang hanya
tersirat pada donor. Suku
$-g_{J_k}^k+\bar g_k$ merupakan peubah kontrol
(\emph{control variate}). Donor juga menyebut SAG dan SVRG; keduanya adalah
metode reduksi varians lain, tetapi algoritmanya tidak diimpor dari bagian di
luar rentang beku.

% B03-S003 | penghubung matematis mandiri untuk baris 2974-2981
% segment-id: d90.becker.98ed693.b03.seg0003
\section{Mengapa koreksi tabel mengurangi varians}
\label{becker03:section:variance-identity}

\noindent\textbf{Penghubung mandiri.}\par
Ambil $\mathcal F_k$ sebagai informasi sebelum $J_k$ ditarik, sehingga
$x_k$ dan seluruh tabel tetap ketika ekspektasi bersyarat dihitung.

\begin{prop}[Takbias bersyarat]
  \label{becker03:prop:conditional-unbiasedness}
  Penaksir~\eqref{becker03:eq:saga-estimator} memenuhi
  \begin{equation}
    \E[v_k\mid\mathcal F_k]=\nabla f(x_k).
  \end{equation}
\end{prop}
\begin{proof}
  Keseragaman $J_k$ memberi
  \begin{multline*}
    \E[v_k\mid\mathcal F_k]
    =\frac1N\sum_{i=1}^{N}
      \bigl(\nabla f_i(x_k)-g_i^k+\bar g_k\bigr)\\
    =\nabla f(x_k)-\bar g_k+\bar g_k
      =\nabla f(x_k).
  \end{multline*}
\end{proof}

Lebih tepat lagi, tuliskan
$a_i^k=\nabla f_i(x_k)-g_i^k$ dan
$\bar a_k=N^{-1}\sum_i a_i^k$. Maka
\begin{equation}
  \label{becker03:eq:variance-identity}
  \E\!\left[
    \norm{v_k-\nabla f(x_k)}_2^2\mid\mathcal F_k
  \right]
  =\frac1N\sum_{i=1}^{N}\norm{a_i^k-\bar a_k}_2^2
  \leq\frac1N\sum_{i=1}^{N}\norm{a_i^k}_2^2.
\end{equation}
Identitas ini memperlihatkan mekanismenya: ketika gradien tersimpan mendekati
gradien komponen pada iterat sekarang, ruas kanan menyusut meskipun langkah
$t_k$ tidak dipaksa menuju nol. Berbeda dari SAG, koreksi SAGA pada
\eqref{becker03:eq:saga-estimator} membuat arah langkah takbias. SVRG mencapai
gagasan serupa dengan gradien penuh pada titik acuan berkala, sehingga tidak
menyimpan seluruh tabel gradien tetapi memakai evaluasi gradien tambahan.

% B03-S004 | APPM5720Notes.tex baris 2982-2988 dengan hipotesis diperbaiki
% segment-id: d90.becker.98ed693.b03.seg0004
\section{Laju konvergensi dan iterat rata-rata}
\label{becker03:section:rates}

Pernyataan donor ``untuk $t$ yang sesuai, metode ini konvergen secara linear''
memerlukan hipotesis. Bentuk berikut adalah spesialisasi nonkomposit dari hasil
SAGA oleh Defazio, Bach, dan Lacoste-Julien (2014).

\begin{theorem}[Laju linear SAGA dengan kekonveksan kuat]
  \label{becker03:theorem:linear-rate}
  Andaikan setiap $f_i$ konveks, terdiferensial, dan mempunyai gradien
  $L$-Lipschitz. Andaikan pula $f$ pada~\eqref{becker03:eq:finite-sum}
  $\mu$-konveks kuat dan mempunyai peminimum $x^*$. Jika seluruh titik tabel
  diinisialisasi pada $x_0$ dan $t_k=t=1/(3L)$, maka
  \begin{equation}
    \label{becker03:eq:linear-rate}
    \E\norm{x_k-x^*}_2^2
    \leq
    \left(1-\min\left\{\frac{1}{4N},\frac{\mu}{3L}\right\}\right)^k C_0,
  \end{equation}
  dengan
  \begin{equation}
    C_0=\norm{x_0-x^*}_2^2
      +\frac{2N}{3L}\bigl(f(x_0)-f(x^*)\bigr).
  \end{equation}
\end{theorem}

Konstanta dan laju pada teorema tersebut mengikuti hasil adaptivitas terhadap
kekonveksan kuat dalam
\href{https://arxiv.org/abs/1407.0202}{Defazio, Bach, dan Lacoste-Julien,
\emph{SAGA} (2014)}. Teorema donor tanpa hipotesis tidak dipertahankan sebagai
klaim universal: tanpa kemulusan dan struktur konveks yang sesuai, laju linear
tidak mengikuti hanya dari bentuk pembaruan.

Jika $f$ hanya konveks, definisikan iterat rata-rata secara tidak ambigu sebagai
\begin{equation}
  \label{becker03:eq:average-iterate}
  \bar x_k=\frac1k\sum_{r=1}^{k}x_r.
\end{equation}
Di bawah asumsi konveks dan $L$-mulus di atas, dengan $t=1/(3L)$, hasil yang
sama memberi batas
\begin{equation}
  \label{becker03:eq:average-rate}
  \E[f(\bar x_k)]-f(x^*)
  \leq\frac{4N}{k}
  \left[
    \frac{2L}{N}\norm{x_0-x^*}_2^2+f(x_0)-f(x^*)
  \right].
\end{equation}
Jadi perataan bukan jaminan tanpa syarat untuk laju yang lebih baik; dalam
rezim konveks yang dinyatakan, ia menyediakan jaminan sublinear
$\mathcal O(1/k)$ untuk nilai fungsi.

% B03-S005 | latihan, petunjuk, dan solusi mandiri
% segment-id: d90.becker.98ed693.b03.seg0005
\section{Latihan dengan petunjuk dan solusi}
\label{becker03:section:exercises}

Latihan berikut merupakan materi baru untuk edisi ini; tidak ada latihan,
petunjuk, atau solusi pada rentang donor.

\begin{exercise}[Takbias dan identitas varians]
  \label{becker03:exercise:variance-identity}
  Dengan $a_i^k=\nabla f_i(x_k)-g_i^k$, buktikan
  Proposisi~\ref{becker03:prop:conditional-unbiasedness} dan
  identitas~\eqref{becker03:eq:variance-identity}. Jelaskan kapan varians
  bersyaratnya sama dengan nol.
\end{exercise}
\noindent\textbf{Petunjuk.}
Gunakan $v_k-\nabla f(x_k)=a_{J_k}^k-\bar a_k$ dan kembangkan kuadratnya.

\noindent\textbf{Solusi lengkap.}
Keseragaman indeks memberi
$\E[a_{J_k}^k\mid\mathcal F_k]=\bar a_k$, sehingga selisih tersebut
bermean nol. Selanjutnya,
\begin{multline*}
  \E\norm{a_{J_k}^k-\bar a_k}_2^2
  =\frac1N\sum_i
    \left(\norm{a_i^k}_2^2-2\inner{a_i^k}{\bar a_k}
      +\norm{\bar a_k}_2^2\right)\\
  =\frac1N\sum_i\norm{a_i^k}_2^2-\norm{\bar a_k}_2^2,
\end{multline*}
yang sama dengan ruas tengah~\eqref{becker03:eq:variance-identity} dan tidak
melebihi ruas kanannya. Variansnya nol tepat ketika semua $a_i^k$ sama;
khususnya, nol jika $g_i^k=\nabla f_i(x_k)$ untuk semua $i$.

\begin{exercise}[Satu langkah kuadratik]
  \label{becker03:exercise:quadratic-step}
  Ambil $N=2$,
  $f_1(x)=\tfrac12(x-1)^2$, $f_2(x)=\tfrac12(x+1)^2$, $x_k=2$,
  $\phi_1^k=0$, dan $\phi_2^k=1$. Hitung dua nilai mungkin bagi $v_k$,
  ekspektasi dan variansnya, lalu bandingkan dengan gradien stokastik biasa
  $\nabla f_{J_k}(x_k)$. Untuk $t_k=0{,}1$, hitung pula dua iterat berikutnya.
\end{exercise}
\noindent\textbf{Petunjuk.}
$g_1^k=-1$, $g_2^k=2$, dan $\bar g_k=1/2$.

\noindent\textbf{Solusi lengkap.}
Pada $x_k=2$, gradien komponen adalah $1$ dan $3$. Oleh karena itu
$v_k=1-(-1)+1/2=5/2$ jika $J_k=1$, sedangkan
$v_k=3-2+1/2=3/2$ jika $J_k=2$. Maka
$\E[v_k]=2=\nabla f(2)$ dan
$\operatorname{Var}(v_k)=\tfrac12[(1/2)^2+(-1/2)^2]=1/4$.
Penaksir SGD biasa bernilai $1$ atau $3$, sehingga mean-nya juga $2$ tetapi
variansnya $1$. Dengan $t_k=0{,}1$, SAGA menghasilkan
$x_{k+1}=1{,}75$ atau $x_{k+1}=1{,}85$.
